\documentclass[UTF8,a4paper,11pt]{ctexart}
\usepackage{geometry}
\geometry{margin=1in}
\usepackage{booktabs}
\usepackage{graphicx}
\usepackage{xcolor}
\usepackage{amsmath}
\usepackage{microtype}
\usepackage{hyperref}
\hypersetup{colorlinks=true,linkcolor=blue!50!black,citecolor=blue!50!black,urlcolor=blue!50!black}

% 中文版 2026-10-12 全文重写，与英文版 PAPER_G3R5.tex（同日修订稿）逐节对应。
\title{三代到巨富：五区三代人生历程与超级富豪的家族出身\\[6pt]\large From Grandparents to Grandees: Three Generations, Five Regions, and the Family Origins of the World's Billionaires}
\author{Ariia Xia\\[2pt]\small 独立研究者}
\date{\today}

\begin{document}
\maketitle

\begin{abstract}
“白手起家”的叙事是否掩盖了亿万富豪背后的家族资本？家庭资源与历史冲击又如何在他们的人生中交织？本文构建了一个跨国、三代的当代超级富豪家族生命历程数据库：101 个亿万富豪家族（707 人），覆盖美国、欧洲、中国、日本、韩国五个地区，按七个角色（本人、父母、祖父母、外祖父母）和九个人生维度编码。本文的分析是探索性、描述性的，得到四个模式。（1）在家庭出身可以判定的 82 位样本富豪中，只有 13.4\% 起步于社会底层；多数人出身于已持有中间层资本的家庭（41.5\%）或延续既有家族财富（35.4\%）。（2）历史冲击只在家庭无力缓冲时才打断教育：1946 年及以前出生的富豪中，低起点或家道中落家庭出身者 12 人中有 8 人基础教育中断，父辈崛起或家业延续家庭出身者 22 人中无一中断（Fisher 精确检验 $p<0.001$）；1946 年以后出生者几乎无人中断。（3）在教育信息已知者中，继承型富豪完成高等教育的比例高于非继承者（92\% 对 57\%），但进入高等教育的比例并不高于家境宽裕的创业者（96\% 对 89\%），差别在于是否读完。（4）记录可见性沿代际、母系和性别急剧下降：101 个家族中只有 2 家记有外婆姓名，同代内姓名可得率的性别差最高达 17 个百分点。队列比较显示，1970 年及以后出生的富豪中技术专业比例更高（59\% 对 37\%，$p=0.056$），但看不出家族出身对他们的作用有所减弱。本文的贡献是：一个开放的七角色数据库；对教育中断的“时代冲击 × 家族缓冲”解释；以及把缺失作为数据的系统论证。

\medskip
\noindent\textbf{关键词：}亿万富豪；社会流动；代际传递；教育；精英再生产；缺失数据
\end{abstract}

% =====================================================================
\section{引言}
\label{sec:intro}

全球超级富豪（福布斯级亿万富豪）在世界人口中不到百万分之一，却掌握着不成比例的经济资源。既有研究多从“财富来源”（白手起家还是继承）刻画这一群体（Kaplan \& Rauh 2013; Lillie, de Alwis \& Tisch 2025; Kendzia, Neville \& Gadgil 2025），或只考察单一世代的截面特征（Lu, Fan \& Fu 2021; Korom, Lutter \& Beckert 2017），很少考察完整的三代人生历程：本人、父母、祖父母与外祖父母的成长环境、教育、职业、行业、婚姻与财富轨迹，如何共同塑造了极端财富的产生。本文构建了一个三代生命历程数据库，覆盖五个地区（美国、欧洲、中国、日本、韩国）的 101 个超级富豪家族、707 名家族成员和 733 条人生变更事件，并以探索性、描述性的设计考察五项指导性预期。

社会流动研究反复表明，家庭背景的作用不止于一代：祖父母（包括外祖父母）的资本经由中间一代传递，会对孙辈的人力资本与职业地位产生独立影响（Anderson, Sheppard \& Monden 2018; Hällsten \& Kolk 2023）。但这类“三代效应”文献几乎都依赖人口登记数据（瑞典、挪威、荷兰等），研究的是普通人口的收入分位，而不是经济金字塔顶端的超级富豪。顶层财富的传递是否遵循同样的代际逻辑？本文把三代生命历程编码范式应用于超级富豪样本，以回应这一空白。

本文的中心问题是：\emph{亿万富豪的人生路径在多大程度上由家庭通道（祖辈资本）决定，又在多大程度上由时代窗口（教育制度变迁、产业革命）决定？}我们从三个方面操作化“家庭通道”：（1）家庭起点的财富—职业结构（父辈崛起／延续／重建／低起点／半延续）；（2）父母与祖父母的教育和职业传递；（3）常被忽视的母系资本来源（外公、外婆）。“时代窗口”则操作化为：（1）教育获得是否受到制度性断裂（战争、革命、户籍制度、经济危机）的影响；（2）个人成长期所处的产业周期（如计算机革命、互联网浪潮）。

我们考察五项指导性预期（第 3 节）：家庭出身（H1）、教育中断（H2）、继承者的教育（H3）、记录可见性（H4）与队列变化（H5）。由于这些预期是在数据收集和编码过程中逐步形成的，本文把分析定位为探索性研究，而不是验证性的假设检验。三个结果最为突出。第一，多数富豪出身于已持有中间层资本或延续既有财富的家庭，真正的低起点者是少数。第二，\textbf{历史冲击只在家庭无力吸收时才打断基础教育}：在 1946 年及以前出生的队列中，教育中断全部集中在低起点和家道中落的家庭，而在之后的队列中几乎消失。第三，\textbf{记录可见性沿代际、母系和性别急剧下降}（只有 2.0\% 的家族记有外婆姓名）。相比之下，两种常见说法得不到支持：继承者的特点在于读完而不是进入高等教育；我们也看不出家族出身对年轻一代富豪的作用有所减弱。

本文有三点贡献。第一，构建了首个“七角色三代 × 九维度 × 五地区”的当代超级富豪跨国生命历程数据库（区别于 Lillie, de Alwis \& Tisch 2025 的截面研究和 O'Brien 2024 的家族网络研究）。第二，揭示了教育中断中的“时代冲击 × 家族缓冲”模式，这是家庭通道与时代窗口交互作用的一个具体、可检验的形式。第三，论证了把缺失作为数据的价值：在检索力度保持一致的前提下，数据可得性的分布本身就能反映记录结构、母系位置与性别不平等。下文结构如下：第 2 节综述文献；第 3 节给出分析框架与指导性预期；第 4 节说明数据与编码；第 5 节报告发现；第 6、7 节为讨论与结论。

% =====================================================================
\section{文献综述与理论定位}
\label{sec:lit}

\subsection{社会流动中的三代效应}

社会分层研究的经典范式关注父子两代的传递，但近来的人口学与社会学文献显示，传递效应可以跨越三代甚至更多代。Anderson, Sheppard \& Monden（2018）系统回顾了 69 项分析，发现其中 58\% 在控制父母特征后，仍观察到祖父母社会经济特征与孙辈教育结果之间的关联；这种关联平均约为父母作用的四分之一，并且在父母资源较少时最强。在台湾等东亚情境中，也有研究发现祖父母对教育获得的正向效应（Chiang \& Park 2015）。Hällsten \& Kolk（2023）把瑞典北部的历史教区记录与现代人口登记相链接，追踪七代人，发现尽管相邻两代之间流动性较高，前代仍在许多代之后持续影响后代的教育、职业与财富。Toft \& Hansen（2022）记录了挪威金融财富前 1\% 群体中的“王朝核心”（另见 Chan \& Boliver 2013 对英国祖父母效应的研究）。\textbf{但这些研究几乎都以普通人口的收入分位或登记数据为对象，没有触及极端上尾——超级富豪。}本文把三代范式向上延伸到超级富豪群体。

\subsection{超级富豪的构成：白手起家与继承}

关于亿万富豪的实证研究近年来不断增加。Lillie, de Alwis \& Tisch（2025）利用 2010--2023 年福布斯全球富豪榜，考察美国亿万富豪的常春藤院校就读情况，发现教育因财富来源和性别不同而承担不同功能，尤其强化了男性继承者的特权。IZA 的一篇工作论文（Kendzia, Neville \& Gadgil 2025）分析了美国福布斯 400 中排名前 100 的白手起家者，强调其成就是“建造”出来的。就中国而言，Lu（2017）追溯了第一代最显赫私营企业家的出身，发现半数以上出生于城市工人或农民等中下阶层家庭，强版本的“政治资本主义”解释得不到多少支持。Lu, Fan \& Fu（2021）利用 2000--2018 年胡润百富榜描绘了中国超级富豪的社会空间，发现其主要由政治资本与文化资本构成。\textbf{这些研究大多是单一世代的截面分析或两代（父母—子女）分析，缺少包括母系在内的“本人—父母—祖父母—外祖父母”完整三代编码。}本文的七角色结构正是与之区别的关键。

\subsection{教育、家庭资本与精英再生产}

教育社会学传统（布迪厄的文化资本理论、柯林斯的文凭主义）强调教育是阶层再生产的机制（Bourdieu 1986; Collins 1979）。在超级富豪中，教育的功能是双重的：它既是人力资本投资（获得能力的真实途径），也是\textbf{地位正当化的符号}（把继承来的财富转化为“应得的”精英身份）。Sato（2025）发现日本父亲失业对子女教育获得有跨代的负面影响，说明家庭断裂对教育的冲击会延续到下一代。Kim（2024）揭示了韩国继承制度改革与性别化代际转移的制度背景。我们要追问的是：在超级富豪样本中，继承者的教育是否特殊到足以说明它承担了\textbf{正当化}功能，还是只反映了富裕家庭普遍存在的阶层再生产。

\subsection{本文的定位：三点贡献}

相对于既有研究，本文的贡献有三：\textbf{（1）结构}——首个“七角色三代结构（本人 + 父母 + 祖父母 + 外祖父母）× 九维度 × 五地区跨国”的当代超级富豪生命历程数据库（O'Brien 2024 的“家族网”研究重建了 1895--1945 年达拉斯上流社会的完整亲属网络，发现其中近 70\% 的人同属一张亲属网，但局限于美国单一城市和特定历史时期，也不涉及当代亿万富豪；Lillie, de Alwis \& Tisch 2025 则是对亿万富豪的截面研究）；\textbf{（2）模式}——揭示教育中断中的“时代冲击 × 家族缓冲”交互（H2），并考察继承者与创业者在教育上的差异（H3）；\textbf{（3）方法}——论证“缺失即数据”，用数据可得性的分布描述记录结构、母系位置与性别不平等（H4）。

% =====================================================================
\section{分析框架与指导性预期}
\label{sec:theory}

\subsection{核心框架：家庭通道与时代窗口}

本文在“\textbf{家庭通道 × 时代窗口}”这一二维框架下组织五项指导性预期：

\begin{itemize}
\item \textbf{家庭通道}：祖辈资本（财富、职业、教育、社会网络）经由父母一代传递给富豪本人的路径，包括父系与母系。通过三个概念操作化：家庭起点结构（H1）；父母／祖辈的教育与职业传递（H2 的一部分）；继承的正当性（H3）。
\item \textbf{时代窗口}：个人成长时所处的制度与产业周期，包括教育制度变迁（战争、革命、户籍、危机）和产业革命（计算机、互联网）。操作化为：教育获得是否受制度性断裂影响（H2）；成长期所处的产业周期（H5）。
\end{itemize}

\textbf{核心命题}：家庭通道与时代窗口是交互关系，而不是此消彼长。家庭决定\emph{缓冲能力与通道}（冲击和失败能否被吸收，能否接触到贵人网络），时代决定\emph{冲击与机会密度}（战争或动荡是否威胁教育，是否存在可利用的产业红利）。这种交互意味着一个可检验的模式：时代层面的冲击主要在家庭缺乏缓冲资源时，才会转化为长期劣势（如教育中断）。教育正处在这个交汇点上，它既是家庭投资的结果，也是时代制度的产物。

\subsection{五项指导性预期}

以下预期是在建库和编码过程中逐步形成的。因此，我们用它们来组织探索性分析，而不是把它们当作预先登记的假设；与预期相悖的模式和相符的模式一并报告。

\begin{itemize}
\item \textbf{H1（家庭出身）}：多数亿万富豪出身于已持有中间层资本（“父辈崛起”）或延续既有家族财富的家庭，而不是来自社会底层。
\item \textbf{H2（时代冲击 × 家族缓冲）}：基础教育中断主要出现在时代层面的冲击（战争、殖民统治、革命、战后动荡）遇上缺乏缓冲资源的家庭之时；在同一时代，父辈崛起或家业延续家庭的子女应很少出现中断，冲击过后出生的队列中也应很少出现。
\item \textbf{H3（继承者的教育）}：继承者比白手起家者更可能完成高等教育，包括海外学历和商科学历。这一模式与“教育为继承地位正当化”的解释相符，但本身不能把正当化与普通的阶层再生产区分开。
\item \textbf{H4（记录可见性）}：个人在公开可检索来源中的可见性随代际距离、沿母系、对女性而下降；缺失的结构本身包含信息。
\item \textbf{H5（队列变化）}：1970 年及以后出生的富豪更集中于技术领域。我们同时考察“计算机革命”叙事所暗示的“家族出身对他们作用减弱”是否成立。
\end{itemize}

% =====================================================================
\section{数据、编码与方法}
\label{sec:data}

\subsection{样本}

样本由从福布斯亿万富豪榜中选取的 \textbf{101 个超级富豪家族}构成，按地区分层：美国 21 个、欧洲 20 个、中国 20 个、日本 20 个、韩国 20 个（东亚合计 60 个，与欧美形成五区对照）。抽样框同时包括现行榜单和历史榜单的成员，以便让已经去世的创始人（如松下幸之助、本田宗一郎、郑周永、乔布斯、坎普拉德）也能作为家族的锚点。在每个地区内，家族是\textbf{按多样性有目的地挑选}的，而不是随机抽取或按财富排名选取：挑选时有意覆盖白手起家者与继承者、不同行业和不同出生年代。因此，这个样本适合比较不同类型的富豪，但各类型在样本中的占比不能用来估计它们在全体亿万富豪中的占比。每个家族编码 \textbf{7 人}：本人（P1）、父亲（P2）、母亲（P3）、祖父（P4）、祖母（P5）、外公（P6）、外婆（P7），即\textbf{完整覆盖父系与母系三代}。某位成员的信息完全无法获得时，其记录仍予保留并编码为“缺失”（见第 \ref{sec:reliability} 节）。全样本共 \textbf{707 人}。本人一代的出生年份从 19 世纪末延续到 1980 年代，既有白手起家者，也有继承者（28 人）。研究采用“一个家族 = 一位富豪”的结构：每个样本单位以一位福布斯级亿万富豪为锚点，围绕其编码三代亲属网络。

\subsection{编码方案：7 人与 9 维度}

每位成员沿 \textbf{9 个维度}编码（CODING\_SPEC v1.3）：

\begin{table}[htbp]
\centering
\caption{九个人生维度的编码内容}
\label{tab:dims}
\small
\begin{tabular}{@{}lp{11cm}@{}}
\toprule
维度 & 内容 \\
\midrule
D1 出生环境 & 出生年份、出生地（城市级，城乡）、出生时家庭经济状况、时代背景 \\
D2 基础教育 & 小学／初中／高中，学校名称、类型、就读年份 \\
D3 高等教育 & 院校、专业、学位、年份、是否完成 \\
D4 家庭结构 & 父母婚姻、出生排行、主要抚养人、重大家庭事件 \\
D5 婚姻状况 & 婚姻次数、配偶背景、结婚／离婚年份 \\
D6 职业轨迹 & 职位—机构—年份序列 \\
D7 行业 & 所在行业及其变化（附年份） \\
D8 财富轨迹 & 分阶段估算（年份、金额、依据） \\
D9 重大人生事件 & 迁移／战争／金融危机／丧亲／贵人／转折点 \\
\bottomrule
\end{tabular}
\end{table}

此外，每个数据格都附有\textbf{可靠度标签}（已证／单源／推断／缺失）和来源 URL；人生\textbf{变更事件}（733 条）另行记录年份、旧值、新值、原因及原因分级（A 本人陈述／B 权威记载／C 研究者推断／D 无法考证）。

\subsection{可靠度与来源分级}
\label{sec:reliability}

全样本 6{,}363 个数据格的可靠度分布见表 \ref{tab:reliability}。\textbf{我们有意把“缺失”作为数据保留，而不是删除}：缺失的分布本身反映了一个社会的“记录可见性”结构（见第 5.4 节）。

\begin{table}[htbp]
\centering
\caption{各亲属角色的可靠度标签（数据格数；括号内为行百分比）}
\label{tab:reliability}
\small
\begin{tabular}{lrrrrr}
\toprule
角色 & 已证 & 单源 & 推断 & 缺失 & 合计 \\
\midrule
本人（P1） & 683 (75.1\%) & 141 (15.5\%) &  53 (5.8\%) &  32 (3.5\%) & 909 \\
父亲（P2） & 316 (34.8\%) & 245 (27.0\%) & 128 (14.1\%) & 220 (24.2\%) & 909 \\
母亲（P3） & 149 (16.4\%) & 236 (26.0\%) & 183 (20.1\%) & 341 (37.5\%) & 909 \\
祖父（P4） & 106 (11.7\%) &  88 (9.7\%)  &  99 (10.9\%) & 616 (67.8\%) & 909 \\
祖母（P5） &  30 (3.3\%)  &  34 (3.7\%)  &  85 (9.4\%)  & 760 (83.6\%) & 909 \\
外公（P6） &  17 (1.9\%)  &  44 (4.8\%)  &  85 (9.4\%)  & 763 (83.9\%) & 909 \\
外婆（P7） &   0 (0.0\%)  &  31 (3.4\%)  &  74 (8.1\%)  & 804 (88.5\%) & 909 \\
\midrule
合计 & 1{,}301 (20.4\%) & 819 (12.9\%) & 707 (11.1\%) & 3{,}536 (55.6\%) & 6{,}363 \\
\bottomrule
\end{tabular}
\end{table}

\textbf{编码一致性检验}：对教育路径分类（高等教育完成／高等教育主动退出／教育剥夺／高中完成但放弃大学／未记载），第一轮编码由 AI 工具生成、作者审核（见第 \ref{sec:ai} 节）。随后作者从中随机抽取 25\% 的子样本（$n=25$，随机种子 20260903），只看 D2／D3 原始文本、不看第一轮编码结果重新编码。一致率为 84.0\%（21/25），Cohen's Kappa = \textbf{0.628}（Landis \& Koch 1977 标准下的“较高一致”）。\textbf{两轮编码并不完全独立}：作者审核了第一轮，又完成了第二轮，所以这一数字应理解为 AI 辅助编码与作者编码之间的一致性检验，而不是独立编码者之间的信度；我们计划请外部人工编码者复核。4 处不一致逐一裁决：一例（欧莱雅第三代继承人）经第三方来源核实，由“未记载”改为“高等教育主动退出”；一例（爱茉莉太平洋继承人）经核实为延世大学学士加康奈尔大学 MBA，由“未记载”改为“完成”；其余两例是编码规范中已明确标注的定义边界案例（Persson 高中未毕业，编为低教育；一位韩国继承人就读的师范学校，编为中等教育层次）。规则编码器的 Kappa 仅为 0.478，说明人工判断明显优于纯规则。

\subsection{分析策略}
\label{sec:strategy}

\begin{itemize}
\item \textbf{描述统计}：各项预期核心分类的分布及地区分解（含跨地区列联表）。当“未知”类别较大时，同时报告以全部 101 例和以已知个案为分母的占比。
\item \textbf{比较}：由于样本较小，我们使用分类计数和 $2\times2$ 比较（双侧 Fisher 精确检验），不使用回归模型。在 $n=101$ 且框架为事后形成的情况下，$p$ 值只是衡量对比强弱的描述性参考，而非验证性检验。
\item \textbf{检索方式}：各角色的信息主要通过通用搜索引擎检索，并借助 AI 进行搜索和整理（见第 \ref{sec:ai} 节）。找不到合格来源时，该数据格编为“缺失”；未设固定的检索时长或正式的停止规则。因此，“缺失”的含义是\emph{在公开可检索来源中查不到}，H4 测量的是在这类来源中的可见性，而不是一般意义上的档案可见性。查到的来源按编码规范（CODING\_SPEC v1.3）的来源层级排序（家谱与档案记录；学术研究；权威传记；相应语言的维基百科；长篇报道；自媒体仅作线索）。编码规范还要求，只出现在中文自媒体中的家庭背景信息，必须由一个独立的非中文来源交叉验证后才能编码。这一规则可以防范广泛流传的谣言，但也会推高中国家族的缺失率。
\item \textbf{缺失处理}：低信息量家族（8 个，填充率 10--33\%，集中于欧洲／日本隐私文化背景和东亚近年白手起家者）保留在开放数据中，并纳入全部分析。
\end{itemize}

\subsection{AI 工具的使用}
\label{sec:ai}

本项目在四个环节使用了生成式 AI 工具：（1）检索和整理公开的人物生平资料，包括用于来源核实的 DataPro 工具；（2）把生平文本编码为九个维度和教育路径分类的第一轮编码，所有编码均经作者审核；（3）论文的起草、编辑和翻译；（4）分析代码的编写。Claude（Anthropic）被用于审阅论文、对照已发表来源核对参考文献，以及修订文本和分析代码。作者已对照所引来源核查 AI 的输出，并对数据和本文全部内容负责。由于 AI 辅助的检索和编码可能遗漏或误读来源，每个编码数据格都附有来源 URL 和可靠度标签，便于独立核查。

\subsection{局限}
\label{sec:limits}

\begin{enumerate}
\item \textbf{有目的抽样}：家族是按类型、行业和年代的多样性挑选的，不是随机抽样。各类别占比（如 H1 的分布、继承者人数）只描述本样本，不应推广到全体亿万富豪；样本内部的对比（如表 \ref{tab:h2x}）受影响较小，但仍以这种挑选为前提。
\item \textbf{记录可得性偏差}：祖辈资料（尤其是外婆）系统性缺失（101 家中只有 2 家记有外婆姓名），样本在结构上偏向“有记录者”。由于检索依赖搜索引擎且没有正式的停止规则，缺失既反映历史记录状况，也反映网络可检索性。这限制了对低起点家庭祖辈的因果推断。
\item \textbf{AI 辅助检索与编码}：第一轮检索和编码由 AI 辅助完成（见第 \ref{sec:ai} 节）。遗漏错误对 H4 尤为重要，因为检索没有找到的来源会被记为缺失。
\item \textbf{样本量}：101 个家族为分层样本，各分类单元格数有限（如低教育仅 9 人），统计检验力有限，结论以模式描述为主。
\item \textbf{财富估算}：D8 财富轨迹为分阶段估算，币种和依据各不相同，跨年代、跨国比较需谨慎。
\item \textbf{自传性偏差}：本人一代的教育／职业信息常来自自传和访谈（A 级原因），存在选择性叙事；我们以交叉来源加以校正。
\item \textbf{因果识别}：本研究是描述性的生命历程研究，不声称识别净因果效应；“家庭通道”与“时代窗口”是解释框架，不是反事实估计。
\item \textbf{没有非富豪对照组}：样本中只有亿万富豪，因此只能描述亿万富豪的家庭背景是什么样，不能说明是哪些因素\emph{造就}了亿万富豪。与总体人口基准率（如按出生队列和国家划分的教育获得）的比较留待未来研究。
\item \textbf{探索性框架}：指导性预期是在与数据反复对照中形成的，结果应视为有待在独立样本中检验的模式。
\item \textbf{类别定义有重叠}：“重建者”和“低起点”由家庭困境来界定，而家庭困境也进入了我们对“断裂”的理解。因此，H2 中这些类别与教育中断之间的关联有一部分是由定义带来的；更有信息量的是同一时代父辈崛起和家业延续家庭中没有出现教育中断这一事实（见第 5.2 节）。
\end{enumerate}

% =====================================================================
\section{发现}
\label{sec:findings}

\subsection{家庭路径：祖辈资本与“断裂—重建”（H1）}

\textbf{分类。}依据家庭起点和父辈轨迹，我们把 101 位富豪分为六类：\textbf{父辈崛起}（34 人，33.7\%）、\textbf{延续}（29 人，28.7\%）、\textbf{未知／数据不足}（19 人，18.8\%）、\textbf{低起点}（11 人，10.9\%）、\textbf{重建}（6 人，5.9\%）和\textbf{半延续}（2 人，2.0\%）（表 \ref{tab:h1}）。“父辈崛起”指富豪本人实现财富跃升时，父辈已在当地积累起相当的资本或地位（如巴菲特的父亲是股票经纪人兼国会议员）。“延续”指财富主要经由家族代际传递（欧洲老钱、东亚财阀）。“重建”指家庭曾拥有资本但遭遇断裂，由富豪本人重新建立（曹德旺、宗庆后、松下幸之助、稻盛和夫、贝索斯、马斯克）。“低起点”严格限定为出生时家庭处于社会底层（乔布斯、奥尔特加、鲁冠球、本田宗一郎、永守重信等）。

\begin{table}[htbp]
\centering
\caption{H1：家庭起点路径分类（$n=101$）}
\label{tab:h1}
\small
\begin{tabular}{lrr}
\toprule
分类 & 人数 & \% \\
\midrule
父辈崛起（父辈已持有中间层资本） & 34 & 33.7 \\
延续（财富跨代传递） & 29 & 28.7 \\
未知／数据不足 & 19 & 18.8 \\
低起点（出生时家庭处于社会底层） & 11 & 10.9 \\
重建（家庭资本断裂，由本人重建） & 6 & 5.9 \\
半延续 & 2 & 2.0 \\
\midrule
合计 & 101 & 100.0 \\
\bottomrule
\end{tabular}
\end{table}

\textbf{核心发现：亿万富豪的崛起既不是“纯粹白手起家”的叙事，也不是“纯粹继承”的叙事。}由于有 19 个家族无法分类，我们同时报告在 82 个出身已知者中的占比：父辈崛起 41.5\%，延续 35.4\%，低起点 13.4\%，重建 7.3\%，半延续 2.4\%。需要说明的是，“父辈崛起”中包括 6 位父亲创办家族企业的二代继承者。如果只看出身已知的 55 位非继承者，也就是通常被冠以“白手起家”标签的群体，其中 28 人（50.9\%）依靠父辈的中间层资本起步，9 人（16.4\%）来自延续的家族财富，只有 11 人（20.0\%）起步于社会底层。可见，即便在创业者中，真正的低起点也是例外。这与对福布斯 400 的研究所记录的家族财富作用“有所下降但依然可观”相一致（Kaplan \& Rauh 2013; Korom, Lutter \& Beckert 2017）。

\textbf{地区差异。}延续型在欧洲（10/20）和韩国（9/20）占比最高，反映老钱和财阀传统；父辈崛起型在美国（11/21）和中国（7/20）占比最高；低起点者分散在各地区（中国 3、日本 3、美国 3、欧洲 2、韩国 0）。“断裂—重建”路径在东亚白手起家者中尤为典型：曹德旺（家中货船沉没）、宗庆后（父辈官宦家庭没落）、稻盛和夫（父亲的纸业作坊衰败）都属于“祖辈资本留有痕迹、由本人重建”的类型。

\subsection{教育中的时代冲击与家族缓冲（H2）}

\textbf{教育是家庭通道与时代窗口交汇得最清楚的地方。}我们把富豪本人的教育路径分为四类：\textbf{高等教育完成}（61 人，60.4\%）、\textbf{高等教育主动退出}（15 人，14.9\%）、\textbf{教育剥夺（基础教育中断）}（9 人，8.9\%）、\textbf{高中完成但放弃大学}（7 人，6.9\%），另有\textbf{未记载}（9 人，8.9\%）。

\textbf{关键口径}：低教育严格界定为最高学历止步于基础教育阶段（小学毕业、中学辍学）；\textbf{没有学位不等于低教育}——进入过大学的人（无论是否毕业）和高中毕业的人都不算低教育。据此，盖茨、扎克伯格、埃里森等“大学退学创业者”归入“高等教育主动退出”，而不是低教育：他们的退学以入学（且多为名校）为前提，是\textbf{主动选择}，与低教育的\textbf{被动剥夺}有本质区别。

\textbf{核心发现 1：9 位低教育富豪中有 8 位处于时代层面的断裂窗口。}八位低教育富豪——奥尔特加（1936，西班牙内战前后）、松下幸之助（1894，明治末期农村破产）、郑周永（1915，日据时期朝鲜贫困山村）、曹德旺（1946，内战—家道中落）、宗庆后（1945，官宦家庭没落）、何享健（1942，战争—贫困）、鲁冠球（1945，战后农村贫困）、本田宗一郎（1906，明治末期铁匠之家）——都出生于 1894--1946 年间本国的战争、殖民统治、革命或战后动荡时期，其中多人还叠加了家庭破产或极度贫困。对他们而言，小学或中学辍学最好理解为时代与社会结构的结果，而非个人选择。\textbf{第九位是明确的例外。}Persson（1979 年生）成长于稳定的瑞典福利国家；他没有读完高中，背景是家庭层面的困境（经济拮据的单亲家庭），而不是时代层面的断裂，并且他很早就转向了编程，所以他的经历中有选择的成分。我们把他视为边界案例：除他之外，样本中没有任何一位富豪在社会稳定时期离开基础教育。

\textbf{核心发现 2：时代冲击只在家庭无力缓冲时才打断教育。}仅有时代层面的冲击并不足以造成教育中断。1946 年及以前出生的 36 位富豪中，有 17 人完成了高等教育，其中许多人与上述八人身处同样饱受战争或战后动荡的社会。表 \ref{tab:h2x} 按出生队列、家庭出身（第 5.1 节的 H1 分类）和基础教育中断情况进行交叉分类。

\begin{table}[htbp]
\centering
\caption{按出生队列和家庭出身划分的基础教育中断情况（$n=101$）}
\label{tab:h2x}
\small
\begin{tabular}{llrrr}
\toprule
出生队列 & 家庭出身（H1） & 人数 & 中断 & \% \\
\midrule
1946 年及以前 & 父辈崛起／延续 & 22 & 0 & 0.0 \\
              & 低起点／重建／半延续 & 12 & 8 & 66.7 \\
              & 未知 & 2 & 0 & 0.0 \\
\midrule
1946 年以后   & 父辈崛起／延续 & 41 & 0 & 0.0 \\
              & 低起点／重建／半延续 & 7 & 0 & 0.0 \\
              & 未知 & 17 & 1 & 5.9 \\
\bottomrule
\end{tabular}

\smallskip
\parbox{0.92\linewidth}{\footnotesize\emph{注}：“中断”指最高学历止步于基础教育阶段（见第 5.2 节）。1946 年以后唯一的中断者是 Persson。双侧 Fisher 精确检验：1946 年及以前队列内的家庭对比，8/12 对 0/22，$p<0.001$；低起点家庭内的队列对比，8/12 对 0/7，$p=0.013$。4 位出生年份字段空缺的富豪按公开资料归入相应队列。}
\end{table}

这张表显示了交互作用的两个方面。\emph{在同一时代内}，家庭出身把结果完全区分开来：1946 年及以前出生、来自父辈崛起或延续家庭的人无一离开基础教育，而来自低起点或家道中落家庭的人有三分之二中断了教育。\emph{在同一类家庭内}，时代同样重要：1946 年以后出生的 7 位低起点富豪中，无人教育中断。由此可见，教育中断出现在时代层面的冲击遇上无力吸收冲击的家庭之时。如第 \ref{sec:limits} 节所述，“重建”和“低起点”的定义部分取决于家庭困境，所以第二行的结果在一定程度上由分类方式决定；第一行的零和队列对比才是更有信息量的结果。Persson 的例子说明家庭困境在任何时代都可能出现，因此未来研究应把时代层面的断裂和家庭层面的断裂分开编码。

\textbf{核心发现 3：主动退学在地区和行业上高度集中。}15 位“高等教育主动退出”者中，有 10 位在科技／互联网行业致富（盖茨、埃里森、扎克伯格、莫斯科维茨、库姆、卡兰尼克、乔布斯、埃克、堀江贵文、普萨），其中 7 位在美国；其余 5 位分属博彩、媒体、时尚、化妆品和人力资源行业。\textbf{主动退学与美国计算机革命的“窗口”紧密耦合}：当技术革命的回报足够高时，大学学位的预期收益会被“时点机会成本”超过。这是“教育—时代”关系的另一面：教育是一道门槛，而“放弃”这道门槛同样取决于时代红利。

\textbf{核心发现 4：完成高等教育者占多数（60.4\%）。}五个地区中，韩国的完成率最高（16/20），中国次之（14/20），反映了东亚对教育投资的高度制度化；美国主动退学者最多（9/21），欧洲“未记载”最多（6/20，与隐私文化有关）。完成高等教育的普遍性，与“亿万富豪的教育程度普遍较高”的既有发现一致（Lillie, de Alwis \& Tisch 2025; Wai 2013）。

\subsection{继承者的教育：差在读完，不在进入（H3）}

样本中有 \textbf{28 位继承型富豪}，即财富来自接掌而非创办家族企业的人：韩国 15 位、欧洲 9 位、日本 3 位、中国 1 位，没有美国继承者。每位继承者的教育都经过人工逐一编码。\footnote{开放数据文件 \texttt{h3\_inheritors\_recode\_v14.csv}。}其中 23 人（82\%）完成了高等教育，1 人进入大学但未完成，1 人完成了中等师范学校，无人基础教育中断，3 人教育情况无记载（韦特海默和劳辛家族的两人）。在教育情况已知的 25 人中，23 人（92\%）获得学位。11 人（39\%）拥有研究生学位，其中 7 人（25\%）为 MBA 或 EMBA；13 人（46\%）曾在成长国以外求学。

\textbf{核心发现：继承者与创业者的区别在于读完高等教育，而不在于进入高等教育。}在教育情况已知的富豪中，继承者获得学位的比例远高于非继承者（92\% 对 57\%；Fisher 精确检验 $p=0.001$）。不过更有说服力的比较对象，是家境宽裕的创业者（非继承的“父辈崛起”者），他们的阶层背景与继承者相近。这些创业者\emph{进入}高等教育的比例与继承者几乎相同（89\% 对 96\%，$p=0.61$），\emph{读完}的比例却低得多（59\% 对 92\%，$p=0.010$）：其中 8 人主动退学。可见，进入高等教育是各类富裕背景共有的；继承者的特点在于读完，并且通常还会加上商科或海外学历（如财阀继承人李在镕、具光谟、辛东彬）。

\textbf{解释。}这一模式与“教育为继承地位正当化”的解释相符：在韦伯的意义上，教育把“继承来的财富”转化为“应得的”精英地位。但它同样符合另外两种更平常的解释：普通的阶层再生产；以及地位稳固的继承者与面临时效性机会的创业者，在读完学位这件事上的机会成本不同。本文的数据无法区分这几种解释。地区也是一个混杂因素：继承者集中在韩国，而韩国富豪整体的学位完成率就很高。要确认教育具有正当化\emph{功能}，需要关于时间点（如是否在接班前不久获得学历）和家族如何公开呈现这些学历的证据。

\textbf{接班前路径。}只有第一轮数据采集中的 9 位继承者（样本扩充到 28 人之前）编码了详细的接班前经历，因此以下描述只是举例，不具代表性。财阀继承人通常要经历 8--12 年的集团内部轮岗（具光谟：LG 电子 → 斯坦福 MBA → 硅谷初创公司员工 → LG 美国法人 → 集团战略部门 → 会长，历时 12 年）。这 9 人无一跨行业转型，只有两人有外部历练（具光谟在硅谷的一年；丰田章男的海外求学和 MBA）。继承者的路径是“教育 + 内部历练”，而不是“跨行业冒险”，这与创业者的路径形成对照。

\subsection{记录可见性：缺失本身就是数据（H4）}

\textbf{核心发现：个人在公开可检索来源中的“可见性”沿代际、母系和性别系统性下降。}我们使用两个指标。\emph{姓名可得率}是某一角色中记有具体姓名的人所占比例；占位性条目（如“姓名未记载”）计为缺失。\emph{维度覆盖率}是九个编码维度中未被标为“缺失”的比例（表 \ref{tab:reliability}）。

\begin{itemize}
\item \textbf{代际递减}（姓名可得率）：本人 100\% → 父亲 69.3\% → 母亲 52.5\% → 祖父 27.7\% → 祖母 10.9\% → 外公 8.9\% → 外婆 2.0\%（101 家中 2 家）。
\item \textbf{母系差距}：母系远不如父系可见——祖父与外公有姓名的家族分别占 27.7\% 和 8.9\%（维度覆盖率分别为 32.2\% 和 16.1\%）。
\item \textbf{性别差距}：在每一代内，女性都比男性更不可见。姓名可得率的性别差：父母一代 16.8 个百分点（69.3\% 对 52.5\%），祖父母一代 16.8 个百分点（27.7\% 对 10.9\%），外祖父母一代 6.9 个百分点（8.9\% 对 2.0\%）；维度覆盖率的性别差分别为 13.3、15.8 和 4.5 个百分点。合并所有非本人角色，男性有姓名的比例为 35.3\%，女性为 21.8\%（相差 13.5 个百分点）。性别差\emph{并不}随代际距离单调扩大：在母系一侧，外公和外婆几乎都没有记录，因此母系位置的影响至少与性别一样大。（以上比较不含本人一代，因为本人以男性为主，且按研究设计全部有记录。）
\item \textbf{地区差异}：韩国财阀家族的祖辈记录最完整（20 个家族中，9 家记有祖父姓名，7 家记有祖母姓名）；中国、日本、韩国的家族中没有一家记有外婆姓名；日本外婆的维度覆盖率最低（2.8\%），这可能与日本《户籍法》对第三方查阅的限制以及母系的边缘化都有关系。
\end{itemize}

\textbf{“缺失即数据”的理论意义。}在检索方式一致的前提下，缺失不只是数据收集的失败，更是\textbf{记录结构的映射}：谁被写进公共记录、进而进入可检索的来源，取决于其在公共领域中的角色、所处的父系或母系位置以及性别。祖母一代不仅缺少姓名，教育、职业和财富信息也几乎全部空白，这是系统性“女性历史不可见”的量化证据。

\textbf{一个重要的限定。}缺失也反映检索本身。信息主要通过搜索引擎收集，查不到即编为缺失，因此缺失在一定程度上测量的是网络可检索性。在编码一致性检验中（第 \ref{sec:reliability} 节），有两例最初编为“缺失”的信息后来被找到；而中文来源的交叉验证规则（第 \ref{sec:strategy} 节）会机械地推高中国家族的缺失率（80 位中国祖辈中只有 3 人有姓名记录）。如果同一家族的各个角色是用同样的来源检索的，家族\emph{内部}各角色之间的差异受这一问题的影响应小于地区之间的差异；因此，我们更看重代际、母系和性别梯度，而不是地区对比。

\subsection{队列变化：年轻一代与年长一代（H5）}

按出生年份把样本分为年轻队列（1970 年及以后出生，$n=34$）和年长队列（1970 年以前出生，$n=67$），同时考察专业领域和家庭出身。

\textbf{发现 1：年轻一代富豪更集中于技术领域。}年轻队列中 59\%（20/34）学的是计算机／工程／物理，年长队列为 37\%（25/67），相差 22 个百分点（Fisher 精确检验 $p=0.056$）。这种集中（张一鸣、黄峥、王兴、扎克伯格、佩奇、程维等）对应 1990 年代末至 2000 年代初的计算机热潮和互联网兴起。由于富豪榜是按财富筛选的，这主要说明近期财富产生于\emph{哪些领域}，并不说明学习技术能提高任何人成为亿万富豪的机会。

\textbf{发现 2：看不出家族出身对年轻一代富豪的作用有所减弱。}如果计算机革命降低了进入门槛，低起点者在年轻队列中应该更多。实际情况正好相反：在出身已知者中，年轻队列只有 1/23（4\%）起步于社会底层，年长队列为 10/59（17\%）（$p=0.17$）；年轻队列的 23 人中有 16 人依靠父辈的中间层资本起步。这一差异在统计上并不可靠，年轻队列 34 人中也有 11 人出身未知，但其方向与“家族作用下降”的叙事相反。个案也说明了同样的问题：王兴（富商家庭，能承受多次创业失败）、张一鸣（工薪家庭，路径线性）、黄峥（工人家庭，靠学霸逆袭和贵人段永平弥补家庭资源）。在每个案例中，家庭提供“容错与通道”，时代提供“机会密度”。

\textbf{地区间的可得性差异本身也是一种发现。}欧洲（继承型财富占主导）和韩国（财阀换代缓慢）的年轻一代白手起家科技富豪都很少——\textbf{年轻一代富豪的国别分布本身反映了各国的创业生态与代际结构。}

% =====================================================================
\section{讨论}
\label{sec:discussion}

\subsection{亿万富豪的形成既非“白手起家”的神话，也非“继承”的宿命}

本文数据呈现的图景比两种流行叙事都更复杂。在出身已知的 82 位富豪中，“父辈崛起”（41.5\%）和“延续”（35.4\%）合计超过四分之三，真正的低起点者占 13.4\%；即便在出身已知的非继承者中，也只有五分之一起步于社会底层。\textbf{对本样本中的多数富豪而言，“白手起家”的财富更准确的描述是“在既有资本之上的攀升”，而不是“从零开始”。}这意味着，亿万富豪研究和大众叙事由于聚焦于白手起家的传奇，可能高估了个人奋斗的分量，低估了家庭通道的作用。这里有两点限定。由于家族是按多样性而非随机挑选的，这些占比描述的是样本而非全体亿万富豪；由于样本中没有非富豪，这些占比也无法估计家族出身能在多大程度上提高成为富豪的几率。

\subsection{教育：时代冲击与家族缓冲的交汇处}

教育方面的发现是本文最有力的部分，它让“家庭通道 × 时代窗口”的框架变得具体。\textbf{单有时代冲击并不会打断教育；冲击遇上无资源的家庭才会。}在 1946 年及以前出生的队列中，战争、殖民统治和战后动荡波及每一个人，但来自父辈崛起或延续家庭的 22 位富豪无一离开基础教育，来自低起点或家道中落家庭的 12 人中却有 8 人中断。在低起点家庭中，这些冲击过去之后，同样的中断就不再出现。1946 年以后唯一的中断者 Persson 说明，在稳定时代，家庭困境也能造成同样的结果。这修正了“富豪的低教育由时代决定”的说法：它由时代决定，\emph{是对无力缓冲时代冲击的家庭而言}。用本文框架的话来说，家庭决定冲击能否被吸收，时代决定冲击是否发生。

继承者和创业者在教育上也有差异，但没有“正当化”叙事所暗示的那么大。家境宽裕的创业者进入高等教育的比例与继承者相当；继承者的不同之处在于\emph{读完}，并加上商科或海外学历。这究竟反映了正当化、普通的阶层再生产，还是地位稳固的继承者读书的机会成本更低，本文的数据无法判断。继承者的路径是“教育 + 内部历练”，而不是跨行业冒险，这与创业者的路径形成对照。

\subsection{“缺失即数据”：记录可见性的结构}

H4 的方法论贡献值得强调。姓名可得率从本人的 100\% 下降到外婆的 2.0\%，母系明显不如父系可见，同代内的性别差最高达 16.8 个百分点——\textbf{在检索方式一致的前提下，缺失映射的是“谁被记录、谁被记住”的社会结构。}祖母一代不仅缺少姓名，教育、职业和财富信息也几乎全部空白，这是系统性“女性历史不可见”的量化证据，也与 Kim（2024）关于韩国性别化代际转移的发现相呼应。

\textbf{一个留待未来研究的猜想：隐私作为一种资源。}8 个低信息量家族（填充率 10--33\%）中，既有欧洲老钱（如匡特、劳辛家族），也有日本近年的富豪。对前者而言，低可见性可能部分源于主动的隐私管理，而这需要财富才能买到的法律资源和媒体关系。如果是这样，可见性就是一个双轨现象：底层家庭是被动地不可见，顶层家庭则可以选择变得不可见。本文的数据无法检验这一点：我们没有衡量隐私管理的指标，而近年白手起家者的低信息量至少同样可能是被动的。检验这一猜想需要专门的“可见性控制”数据，如传记授权记录或针对媒体的法律行动记录。

\subsection{队列变化：新产业，旧优势}

H5 显示，年轻一代富豪中有 59\% 学的是计算机、工程或物理，比年长一代高 22 个百分点。技术变迁改变了通向极端财富的\emph{领域}，却没有为缺乏家庭资源的人把门开得更大：在出身已知的年轻一代富豪中，低起点者更少而不是更多，他们中的多数依靠父辈的中间层资本起步。“计算机革命削弱了家庭背景的作用”这一常见说法在这里得不到支持。如果说有什么不同，一个奖励早期技术教育、需要承受多次创业失败的行业，可能更有利于能同时提供这两者的家庭。王兴（富裕家庭，能承受多次失败）和黄峥（工人家庭，靠学霸逆袭和贵人段永平弥补资源不足）代表了两条路径：\textbf{家庭提供“容错与通道”，时代提供“机会密度”。}

\subsection{与文献的对话}

与既有文献相比，本研究利用三代、含母系、五地区的数据，把亿万富豪研究从“截面画像”推向“代际过程”。本文的发现与两类文献相关联。对教育不平等研究而言，表 \ref{tab:h2x} 中的“时代 × 家庭”模式呼应了家庭冲击会降低子女教育的发现（Sato 2025）：在 1946 年及以前出生的队列中，承担历史冲击教育代价的是没有缓冲的家庭的孩子，而不是父辈崛起或家业延续家庭的孩子。对精英再生产研究而言（Bourdieu 1986; Collins 1979; Lillie, de Alwis \& Tisch 2025），继承者高等教育完成率高、且多有精英商科学历，符合“教育帮助把继承优势转化为被认可的功绩”的论述，尽管本文数据无法把这一机制与普通的阶层再生产区分开。从 101 位亿万富豪的样本中引出政策含义需要谨慎。最稳妥的启示是一个一般性的结论：历史冲击会不会变成长期的教育劣势，很大程度上取决于家庭能否吸收它。

% =====================================================================
\section{结论}
\label{sec:conclusion}

本研究构建了一个“七角色三代结构”的生命历程数据库，覆盖美国、欧洲、中国、日本和韩国的 101 个超级富豪家族、707 名成员和 733 条人生变更事件，并以探索性设计考察了五项指导性预期。主要结论如下。

\textbf{第一，本样本中的多数亿万富豪并非从零起步。}在出身已知者中，超过四分之三依靠父辈的中间层资本起步（41.5\%）或延续既有家族财富（35.4\%），只有 13.4\% 起步于社会底层；即便在创业者中，真正的低起点也是例外。

\textbf{第二，家庭通道与时代窗口相互作用，而教育是这种交互最清楚的体现。}历史冲击只在家庭无力吸收时才打断基础教育：1946 年及以前出生的富豪中，低起点或家道中落家庭出身者 12 人中有 8 人离开基础教育，父辈崛起或延续家庭出身者 22 人中无一人如此；在之后的队列中，这种中断几乎消失。这一“时代冲击 × 家族缓冲”模式是本文最主要的实质性发现。

\textbf{第三，继承者的特点在于读完高等教育，而不在于进入高等教育。}继承者获得学位的比例远高于家境宽裕的创业者（92\% 对 59\%），而两者进入高等教育的比例相近。这与教育承担正当化功能的解释相符，但不足以证明这一点。

\textbf{第四，记录可见性具有高度的结构性。}姓名可得率从本人的 100\% 下降到外婆的 2.0\%，母系远不如父系可见，同代内的性别差最高达 16.8 个百分点。在检索力度保持一致的前提下，缺失本身就是关于“谁被记录、谁被记住”的证据。

\textbf{第五，新产业并没有降低旧门槛。}年轻一代富豪更集中于技术领域（59\% 对 37\%），但其中低起点者更少，而不是更多。

这些结论是描述性、探索性的。样本规模小且为有目的挑选，没有非富豪对照组，依赖 AI 辅助的网络检索和编码，编码框架也是与数据一同形成的。尤其是“时代冲击 × 家族缓冲”模式，应当在独立样本中、并对照按队列和国家划分的总体教育基准率加以检验。未来研究还可以把时代层面与家庭层面的断裂分开编码，收集关于继承者学历获得时间及其公开呈现方式的证据，并用专门数据发展“隐私作为资源”的猜想。我们公开全部数据（含编码规范和字段字典），以支持这一新兴领域的可复现研究。

\textbf{总之}：亿万富豪的诞生既不是纯粹的个人神话，也不是简单的继承宿命。家庭资源决定历史冲击能否被吸收、机会能否被抓住；时代决定哪些冲击和机会会到来。理解这种交汇，不仅关乎对极富群体的刻画，也关乎对机会平等和阶层再生产的理解。

% =====================================================================
\begin{table}[htbp]
\centering
\caption{按原因分级划分的人生变更事件（$n=733$）}
\label{tab:events}
\small
\begin{tabular}{lrr}
\toprule
分级 & 条数 & \% \\
\midrule
A（本人陈述） & 124 & 16.9 \\
B（权威记载） & 451 & 61.5 \\
C（研究者推断） & 153 & 20.9 \\
D（无法考证） & 5 & 0.7 \\
\midrule
合计 & 733 & 100.0 \\
\bottomrule
\end{tabular}
\end{table}

% =====================================================================
\section*{数据可得性}
完整的编码数据库（707 人 × 9 维度，附可靠度标签和来源 URL）、事件记录（733 条变更）、编码规范（CODING\_SPEC v1.3）和教育路径分类，均作为本文的开放数据公开。

% =====================================================================
\begin{thebibliography}{99}
\setlength{\itemsep}{2pt}

\bibitem{anderson2018} Anderson, L. R., Sheppard, P., \& Monden, C. W. S. (2018). Grandparent effects on educational outcomes: A systematic review. \emph{Sociological Science}, 5, 114--142.

\bibitem{bourdieu1986} Bourdieu, P. (1986). The forms of capital. In J. G. Richardson (Ed.), \emph{Handbook of Theory and Research for the Sociology of Education} (pp. 241--258). Greenwood.

\bibitem{chan2013} Chan, T. W., \& Boliver, V. (2013). The grandparents effect in social mobility: Evidence from British birth cohort studies. \emph{American Sociological Review}, 78(4), 662--678.

\bibitem{chiang2015} Chiang, Y.-L., \& Park, H. (2015). Do grandparents matter? A multigenerational perspective on educational attainment in Taiwan. \emph{Social Science Research}, 51, 163--173.

\bibitem{collins1979} Collins, R. (1979). \emph{The Credential Society: An Historical Sociology of Education and Stratification}. Academic Press.

\bibitem{hallsten2023} Hällsten, M., \& Kolk, M. (2023). The shadow of peasant past: Seven generations of inequality persistence in Northern Sweden. \emph{American Journal of Sociology}, 128(6), 1716--1760.

\bibitem{kaplan2013} Kaplan, S. N., \& Rauh, J. D. (2013). Family, education, and sources of wealth among the richest Americans, 1982--2012. \emph{American Economic Review}, 103(3), 158--162.

\bibitem{iza2026} Kendzia, M. J., Neville, T., \& Gadgil, M. (2025). \emph{Built, not born: How education predicts billionaire wealth}. IZA Discussion Paper No. 17994. IZA Institute of Labor Economics.

\bibitem{kim2025} Kim, D. (2024). An equal right to inherit? Inheritance rights and gendered intergenerational transfers in South Korea, 1971--2010. \emph{Population}, 79(2--3).

\bibitem{korom2017} Korom, P., Lutter, M., \& Beckert, J. (2017). The enduring importance of family wealth: Evidence from the Forbes 400, 1982 to 2013. \emph{Social Science Research}, 65, 75--95.

\bibitem{landis1977} Landis, J. R., \& Koch, G. G. (1977). The measurement of observer agreement for categorical data. \emph{Biometrics}, 33(1), 159--174.

\bibitem{lillie2025} Lillie, K., de Alwis, C., \& Tisch, D. (2025). Male heirs and self-made women? Gender, extreme wealth and the Ivy League. \emph{AERA Open}, 11.

\bibitem{lu2017} Lu, P. (2017). The Horatio Alger myth in China: Origins of the first generation of visibly richest Chinese private entrepreneurs. \emph{China: An International Journal}, 15(2), 75--97.

\bibitem{lu2021} Lu, P., Fan, X., \& Fu, F. (2021). Profile of the super rich in China: A social space analysis. \emph{British Journal of Sociology}, 72(3), 543--565.

\bibitem{familyweb2023} O'Brien, S. (2024). The family web: Multigenerational class persistence in elite populations. \emph{Socio-Economic Review}, 22(1), 1--27.

\bibitem{sato2024} Sato, K. (2025). Intergenerational associations between paternal job loss and children's educational attainment in Japan. \emph{Journal of the Japanese and International Economies}, 75, 101341.

\bibitem{toft2022} Toft, M., \& Hansen, M. N. (2022). Dynastic cores and the borrowed time of newcomers: Wealth accumulation and the Norwegian one percent. \emph{British Journal of Sociology}, 73(2), 291--314.

\bibitem{wai2013} Wai, J. (2013). Investigating America's elite: Cognitive ability, education, and sex differences. \emph{Intelligence}, 41(4), 203--211.

\end{thebibliography}

\end{document}
